% Part: proof-theory
% Chapter: natural-deduction
% Section: rules-proofs

\documentclass[../../../include/open-logic-section]{subfiles}

\begin{document}

\olfileid[pt]{pt}{nat}{rp}

\olsection{Regras e \usetoken{P}{proof}}

Começamos apresentando algumas definições e fixando certa terminologia.

Considere as duas regras de \Log{N1c} e~\Log{N1i} que mencionam $!A
\lif !B$.
\[
\bottomAlignProof
\AxiomC{$\Discharge{!A}{x}$}
\shortDeduce
\DeduceC{$!B$}
\DischargeRule{\Intro{\lif}}{x}
\UnaryInfC{$!A \lif !B$}
\DisplayProof
\qquad
\bottomAlignProof
\AxiomC{$!A \lif !B$}
\AxiomC{$!A$}
\RightLabel{\Elim{\lif}}
\BinaryInfC{$!B$}
\DisplayProof
\]
A regra \Elim{\lif} é bastante simples. As !!{formula}s acima da linha
são as \emph{premissas}. A !!{formula} à esquerda é a
!!{formula} \emph{principal}~$!A \lif !B$. Todas as regras de !!{elimination}
têm uma premissa desse tipo, sempre na extrema esquerda. Ela é chamada de
premissa \emph{maior} da inferência. Se há outras, como neste
caso, elas são chamadas de premissas \emph{menores}. As premissas das
regras de !!{introduction} também são chamadas de premissas maiores. A
!!{formula} abaixo da linha é a \emph{conclusão}.

A regra \Intro{\lif} tem apenas uma premissa, e aqui a conclusão
é a fórmula principal~$!A \lif !B$. Todas as regras de !!{introduction} são
assim: a conclusão é a fórmula principal; seu operador principal
é o que está sendo ``introduzido''.

A notação $\Discharge{!A}{x}$ significa o seguinte. Em !!a{proof},
a inferência \Intro{\lif} é aplicada a !!a{formula}~$!B$ como premissa.
A sub-!!{proof} que termina em~$!B$ tem várias !!{formula}s iniciais.
Elas são chamadas de \emph{hipóteses}. !!^a{proof} de~$!B$ a partir de uma ou
mais hipóteses da forma~$!A$ é !!a{proof} de que $!A$
implica~$!B$. Quando aplicamos a regra \Intro{\lif}, concluímos~$!A
\lif !B$, e a conclusão já não depende da hipótese~$!A$.
Para indicar isso, as hipóteses da forma~$!A$ na !!{proof}
que inclui a inferência~\Intro{\lif} são \emph{descarregadas} ou
\emph{fechadas}. A regra especifica quais hipóteses podem ser descarregadas. No
caso de \Intro{\lif} com conclusão $!A \lif !B$, são as
hipóteses que têm a forma~$!A$, isto é, que coincidem com o antecedente
da condicional. Às vezes é importante indicar quais
hipóteses são descarregadas em qual regra, e isso é feito
com o rótulo de descarga~$x$. É importante notar que as regras que descarregam
hipóteses não \emph{precisam} descarregar \emph{todas} as hipóteses da
forma exigida, nem sequer alguma delas. Por exemplo, a seguinte é
correta, embora a primeira \Intro{\lif} não descarregue nenhuma
hipótese:
\[
\AxiomC{$\Discharge{!A}{y}$}
\DischargeRule{\Intro{\lif}}{x}
\UnaryInfC{$!B \lif !A$}
\DischargeRule{\Intro{\lif}}{y}
\UnaryInfC{$!A \lif (!B \lif !A)$}
\DisplayProof
\]
Quando uma regra não descarrega nenhuma hipótese, simplesmente omitimos o
rótulo de descarga (como fizemos aqui).

Na !!{proof}
\[
\AxiomC{$\Discharge{!A}{x}$}
\AxiomC{$\Discharge{!A}{y}$}
\RightLabel{\Intro{\land}}
\BinaryInfC{$!A \land !A$}
\RightLabel{\Elim{\land}}
\UnaryInfC{$!A$}
\DischargeRule{\Intro{\lif}}{x}
\UnaryInfC{$!A \lif !A$}
\DischargeRule{\Intro{\lif}}{y}
\UnaryInfC{$!A \lif (!A \lif !A)$}
\DisplayProof
\]
a primeira inferência \Intro{\lif} descarrega a hipótese à esquerda
$!A^x$, e a segunda, a hipótese à direita~$!A^y$. Isto é, todas as
hipóteses rotuladas por~$x$ são descarregadas pela \Intro{\lif} com
rótulo~$x$, e todas as hipóteses com rótulo~$y$, por aquela rotulada
por~$y$. Para que isso funcione, é importante que todas as hipóteses
com o mesmo rótulo sejam também a mesma !!{formula}.

As regras dos quantificadores são um pouco mais complicadas. Por
exemplo, as regras de \Log{N1i} para $\lexists$ são:
\[
\bottomAlignProof
\AxiomC{$!A(t)$}
\RightLabel{\Intro{\lexists}}
\UnaryInfC{$\lexists[x][!A(x)]$}
\DisplayProof
\qquad
\bottomAlignProof
\AxiomC{$\lexists[x][!A(x)]$} 
\AxiomC{$\Discharge{!A(c)}{x}$}
\shortDeduce
\DeduceC{$!B$}
\DischargeRule{\Elim{\lexists}}{x}
\BinaryInfC{$!B$} \DisplayProof 
\bottomAlignProof
\] Na regra
\Intro{\lexists}, a premissa é~$!A(t)$, em que $t$~é um termo fechado,
isto é, um termo sem variáveis. Uma inferência desse tipo é
correta---isto é, corresponde à regra esquemática \Intro{\lexists}---se
existe alguma !!{formula}~$!A$, alguma variável~$x$ e algum termo
fechado~$t$ tais que a conclusão é $\lexists[x][!A]$ e a premissa
é~$\Subst{!A}{t}{x}$. A regra \Elim{\lexists} permite descarregar
hipóteses da forma~$!A(c)$, isto é, para cada inferência existe alguma
!!{constant}~$c$, e hipóteses da forma $\Subst{!A}{c}{x}$ podem ser
descarregadas. Todas as hipóteses descarregadas por uma inferência
devem usar o mesmo~$c$. Essa constante é chamada de
\emph{variável própria} da inferência \Elim{\lexists}. A inferência só
é correta se $c$~não ocorre em nenhuma hipótese aberta além de
$!A(c)$, que é descarregada, não ocorre em~$!B$ e não ocorre
em~$!A$. Isso é chamado de \emph{condição da variável própria}.

!!^{proof}s em \Log{N1c} e~\Log{N1i} são árvores finitas de sequentes
geradas indutivamente a partir de axiomas mediante as regras de
inferência. Cada folha é uma hipótese, possivelmente marcada com um
rótulo de descarga, e toda outra !!{formula} segue das
!!{formula}s imediatamente acima dela por uma das regras de inferência do
sistema. Se uma ocorrência de uma !!{formula} como hipótese na árvore
acima de uma inferência não foi descarregada por uma inferência acima dela,
ela é chamada de \emph{aberta} (na inferência).

\begin{defn}\ollabel{defn:proof}
\emph{!!^{proof}s} em \Log{N1c} e~\Log{N1i} são árvores finitas de
!!{formula}s definidas indutivamente da seguinte maneira:
\begin{enumerate}
  \item\ollabel{defn:prf:assumption} Toda !!{formula}
  marcada com um rótulo de descarga, como $x$ ou $y$, é, por si só,
  !!a{proof}. Em !!a{proof} que consiste em uma única !!{formula}, a
  !!{formula} conta como hipótese aberta.
  \item\ollabel{defn:prf:one-prem} Se $R$ é uma regra com uma, duas ou
  três premissas $!A_1$, $!A_2$, $!A_3$, conclusão~$!A$, e
  $\delta_1$, $\delta_2$, $\delta_3$ são !!{proof}s que terminam em $!A_1$,
  $!A_2$, $!A_3$, respectivamente, então cada uma das construções
  \[
  \AxiomC{}
  \RightLabel{$\delta_1$}
  \DeduceC{$!A_1$}
  \RightLabel{$R$}
  \UnaryInfC{$!A$}
  \DisplayProof
\qquad
  \AxiomC{}
  \RightLabel{$\delta_1$}
  \DeduceC{$!A_1$}
  \AxiomC{}
  \RightLabel{$\delta_2$}
  \DeduceC{$!A_2$}
  \RightLabel{$R$}
  \BinaryInfC{$!A$}
  \DisplayProof
\qquad
  \AxiomC{}
  \RightLabel{$\delta_1$}
  \DeduceC{$!A_1$}
  \AxiomC{}
  \RightLabel{$\delta_2$}
  \DeduceC{$!A_2$}
  \AxiomC{}
  \RightLabel{$\delta_3$}
  \DeduceC{$!A_2$}
  \RightLabel{$R$}
  \TrinaryInfC{$!A$}
  \DisplayProof
  \]
  é !!a{proof}, conforme o caso, desde que os rótulos das hipóteses nas
  !!{proof}s sejam compatíveis (nenhuma dupla de !!{formula}s diferentes
  tem o mesmo rótulo de descarga) e as inferências sejam aplicações
  corretas de~$R$.
  \item As ocorrências mais altas de !!{formula} na nova !!{proof} contam
  como \emph{hipóteses abertas} da !!{proof} se são hipóteses abertas de
  $\delta_1$, $\delta_2$ ou~$\delta_3$, exceto pelas ocorrências
  descarregadas pela regra. Se a regra tem o rótulo~$x$ (ou $x\, y$),
  então, para \Intro{\lif} e~\Intro{\lnot}, essas ocorrências são todas
  as hipóteses abertas de~$\delta_1$ rotuladas por~$x$; para
  \Elim{\lexists}, são todas as hipóteses abertas rotuladas por~$x$ em
  $\delta_2$; para \Elim{\lor}, são todas as hipóteses abertas rotuladas
  por~$x$ em~$\delta_2$ e todas as hipóteses abertas rotuladas por~$y$
  em~$\delta_3$.
\end{enumerate}
\end{defn}

As !!{proof}s usadas na construção indutiva de !!a{proof} são chamadas
de suas \emph{sub-!!{proof}s}. As sub-!!{proof}s que terminam nas
premissas de uma inferência são chamadas de \emph{sub-!!{proof}s
imediatas} da inferência. As regras de \Log{N2c} e \Log{N2i} são
apresentadas em \olref{tab:N1}.

\subfile{rules-N1}

\begin{defn}
A \emph{!!{height}}~$\pheight{\delta}$ de !!a{proof}~$\delta$ é
definida indutivamente da seguinte maneira.
\begin{enumerate}
  \item Se $\delta$ consiste em uma hipótese, $\pheight{\delta} = 0$.
  \item Se $\delta$ termina em uma inferência com uma premissa cuja
  sub-!!{proof} imediata é~$\delta_1$, então $\pheight{\delta} = \pheight{\delta_1} + 1$.
  \item Se $\delta$ termina em uma inferência com duas premissas cujas
  sub-!!{proof}s imediatas são~$\delta_1$ e~$\delta_2$, então $\pheight{\delta} =
  \max(\pheight{\delta_1}, \pheight{\delta_2}) + 1$.
  \item Se $\delta$ termina em \Elim{\lor} com sub-!!{proof}s imediatas
  $\delta_1$, $\delta_2$ e~$\delta_3$, então $\pheight{\delta} =
  \max(\pheight{\delta_1}, \pheight{\delta_2}, \pheight{\delta_3}) + 1$.
\end{enumerate}
Se um sequente~$S$ tem !!a{proof} de altura $\le n$, escrevemos $\Proves[n] S$.
\end{defn}

\begin{ex}
Em \Log{N1i}, toda !!{formula} conta como !!a{proof} a partir de si
mesma como hipótese aberta. Assim,
\[
!A \land !B^x
\]
é !!a{proof}~$\delta_1$. Sua altura é~$0$.
A partir de $\delta_1$, podemos obter duas !!{proof}s~$\delta_2$ e~$\delta_3$
aplicando uma das duas versões de~\Elim{\land}:
\[
\AxiomC{$!A \land !B^x$}
\RightLabel{\Elim{\land}[2]}
\UnaryInfC{$!B$}
\DisplayProof
\qquad
\AxiomC{$!A \land !B^x$}
\RightLabel{\Elim{\land}[1]}
\UnaryInfC{$!A$}
\DisplayProof
\]
As sub-!!{proof}s imediatas da última inferência em~$\delta_2$ e
em~$\delta_3$ são iguais: $!A \land !B$. Ambas as !!{proof}s têm a
ocorrência de $!A\land !B$ como hipótese aberta. Temos
$\pheight{\delta_1} = \pheight{\delta_2} = 1$.

A regra \Intro{\land} exige duas premissas da forma $!B$ e $!A$
para justificar uma conclusão da forma~$!B \land !A$. Assim, a seguinte
é !!a{proof}~$\delta_4$:
\[
\AxiomC{$!A \land !B^x$}
\RightLabel{\Elim{\land}[2]}
\UnaryInfC{$!B$}
\LeftSubproofLabel{$\delta_2$}
\AxiomC{$!A \land !B^x$}
\RightLabel{\Elim{\land}[1]}
\UnaryInfC{$!A$}
\RightSubproofLabel{$\delta_3$}
\RightLabel{\Intro{\land}}
\BinaryInfC{$!B \land !A$}
\DisplayProof
\]
Sua altura é $\pheight{\delta_4} = \max(\pheight{\delta_2},
\pheight{\delta_3})+1 = \max(1,1) + 1 = 2$. Ela tem duas ocorrências de
$!A \land !B$ como hipóteses, ambas abertas.

Por fim, podemos aplicar \Intro{\lif} para obter $\delta_5$:
\[
\AxiomC{$\Discharge{!A \land !B}{x}$}
\RightLabel{\Elim{\land}[2]}
\UnaryInfC{$!B$}
\LeftSubproofLabel{$\delta_2$}
\AxiomC{$\Discharge{!A \land !B}{x}$}
\RightLabel{\Elim{\land}[1]}
\UnaryInfC{$!A$}
\RightSubproofLabel{$\delta_3$}
\RightLabel{\Intro{\land}}
\BinaryInfC{$!B \land !A$}
\DischargeRule{\Intro{\lif}}{x}
\UnaryInfC{$(!A \land !B) \lif (!B \land !A)$}
\DisplayProof
\]
Sua altura é $\pheight{\delta_4} + 1 = 3$. A inferência \Intro{\lif}
descarrega todas as ocorrências de hipóteses da forma $!A \land !B$
rotuladas por~$x$, isto é, ambas as hipóteses anteriormente abertas da
subprova imediata. Como essas são as únicas hipóteses de $\delta_5$, ela
não tem hipóteses abertas e, portanto, é !!a{proof} \emph{de} sua
!!{formula} final $(!A \land !B) \lif (!B \land !A)$.
\end{ex}

\end{document}
